\chapter{Toán Vô Tội}
\chapterintro{Nhiều khi môn Toán bị ghét oan chỉ vì nó xuất hiện đúng chỗ người ta chưa chịu học}{Không phải lỗi nào của học sinh cũng nên đổ hết cho con số.}

\begin{lucbatbox}{Lục bát: Toán chưa kịp làm gì}
Toán kia nằm đó hiền ghê\
Chỉ ghi vài dấu, vài đề, vài trang\
Em nhìn đã vội than vang\
Như môn kia đã sẵn sàng hại ai\
Nghĩ đi cũng thấy thương thay\
Toán chưa có tội đã ngay án rồi
\end{lucbatbox}

\begin{tutuyetbox}{Thất ngôn tứ tuyệt: Oan cho phương trình}
Phương trình đứng giữa trang im\
Không la, không mắng, không tìm gây chi\
Chỉ vì em ngại nghĩ suy\
Nên môn vô tội cũng vì thành oan
\end{tutuyetbox}

\begin{ghichu}
Toán vô tội ở chỗ nó chỉ đòi người học làm thật. Phần đáng sợ nhiều khi nằm ở thói quen tránh né của mình hơn là ở bài toán.
\end{ghichu}

\ngancach

\begin{lucbatbox}{Lục bát: Ghét oan vì chưa quen}
Mới nhìn đã thấy không ưa\
Vì bao lần trước mình chưa hiểu gì\
Nhưng mà trách hết làm chi\
Khi tay chưa chịu ngồi ghi một hàng\
Muốn thôi cảm giác hoang mang\
Cũng nên thử bước dịu dàng một phen
\end{lucbatbox}

\begin{tutuyetbox}{Thất ngôn tứ tuyệt: Con số bị mang tiếng}
Con số nào biết nói đâu\
Mà nghe bao tiếng trách sầu gọi tên\
Nếu em chịu học từng đêm\
Biết đâu môn ấy lại mềm hơn xưa
\end{tutuyetbox}